\chapter{Netz, Materialien und Modellaufbau}\label{sec:app-models}
\index{Model}\index{Randbedingungen}\index{Elektrode}
\section{Netze erstellen oder importieren}
Direkte Python-Rezepte bauen \path{Mesh(points=..., tetra4=...)} aus SI-Koordinaten
und nullbasierten Knotenverbindungen. Die stationären Quaderbeispiele zeigen
die vollständige Erzeugung. Ein \path{Mesh} kann Tet4 oder geradkantige Tet10
und die zugehörigen Triangle3-/Triangle6-Randflächen enthalten. Volumenzell-
und Flächensätze sind benannte Auswahlmengen, keine Materialobjekte.
Ungenutzte Knoten tragen keine Feldfreiheitsgrade. Tet10-Kantenknoten müssen
mittig auf geraden Kanten liegen; gekrümmte isoparametrische Elemente gehören
nicht zu diesem Funktionsstand. Geometrie- und Orientierungskonventionen
stehen im Theorieband.

Gmsh 2.2 und 4.1 werden über \path{read_gmsh} importiert. Für eine neue
Geometrie kann Gmsh mit \path{gmsh body.geo -3 -format msh2 -o body.msh}
ein lineares Netz erzeugen. Volume- und Surface-Physical-Groups müssen im
\path{.geo} vorgegeben werden; deren Namen werden in TOML referenziert.
Das ist ein optionaler Vorbereitungsschritt, kein mit dem Wheel geliefertes
Meshing-Programm. Für die eingebetteten Rezepte sind alle \path{.msh}-Dateien
bereits enthalten. Bei eigener Tet10-Erzeugung vor dem Solve geradkantige
mittige Kantenknoten und vollständige Triangle6-Ränder kontrollieren.

\path{read_gmsh(path, length_scale_to_m=...)} verlangt eine explizite positive
Längenskala. Ein Netz in mm verwendet $10^{-3}$; Materialwerte, Lasten und
Anschlussflächen bleiben dennoch SI. Die Skala verändert alle geometrischen
Volumina und Flächen. Bei einem fehlerhaft als m interpretierten mm-Netz
korrigiert daher kein nachträglicher Faktor am Verschiebungsplot die Rechnung.

\section{Direkte Python-API oder deklaratives Modell}
Die direkten Solver akzeptieren Materialobjekte, Mesh, DOF-basierte
Dirichlet-Bedingungen und optional räumliche Callables. Beispiele wie
\ref{recipe:expansion} und \ref{recipe:tet10-thermomechanics} verwenden
\path{solve_thermomechanics} direkt. Für Studien und wiederholbare
Dateiverträge verwenden die übrigen Rezepte \path{Model.from_file}.
\begin{lstlisting}
from fem_toolbox.model import Model
model = Model.from_file("examples/thermopiezoelectric/case.toml")
run = model.solve()
result = run.field_result
print(result.heat.temperature)
print(result.piezoelectric.potential)
\end{lstlisting}
Dieses Kurzbeispiel ist ein Ausschnitt aus dem vollständigen eingebetteten
Programm \ref{recipe:thermopiezo}. \path{to_document()} liefert eine unabhängige
Dokumentkopie. Nach bewussten Änderungen konstruiert
\path{Model(mesh=model.mesh, document=document)} ein neues Modell; es validiert
erneut. Der relative Netzpfad wird beim Dateiladen aus dem Modellverzeichnis
aufgelöst. Dateien, die fremde Eingaben referenzieren, werden nicht zu
beliebigen Python-Programmen ausgewertet.

\section{Drei explizite Modellschemata}
\begin{longtable}{p{.12\linewidth}p{.32\linewidth}p{.46\linewidth}}
\toprule Schema & Physik & Dokumentstruktur\\\midrule\endhead
1 & \path{heat}, \path{elasticity}, \path{electrostatics},
\path{piezoelectric}, \path{thermopiezoelectric}
& \path{units="SI"}, Netz, feldweise Materialien/Regionen,
\path{boundary}, \path{loads}; stationär.\\
2 & \path{transient_heat}
& Netz, Materialkapazität, \path{regions}, Zeitachse, Anfangstemperatur,
Zeitprofile, Probes und Ausgabesteuerung.\\
3 & \path{piezo_modes}, \path{transient_piezoelectric},
\path{harmonic_piezoelectric}
& Netz und Dichte, eingebetteter stationärer Schema-1-Piezoteil
\path{base}, analyseabhängige \path{analysis}.\\\bottomrule
\caption{Modellschemata sind Physikverträge, keine Paketversionsnummern.}
\end{longtable}
Schema 1 ist nicht durch 2 oder 3 ersetzt. Alte stationäre Dateien behalten
ihre Bedeutung. Eine andere Zahl einzutragen migriert weder einen
stationären Wärmelauf in einen transienten Lauf noch eine statische
Piezorechnung in Dynamik. Unbekannte Schlüssel, Einheiten oder Versionen
werden abgelehnt. Die frühen Vertragsdokumente zu Schema 1 sind durch die
späteren Schema-2/-3- und Tet10-Erweiterungen ergänzt; diese Übersicht beschreibt
den aktuellen Gesamtstand. Es gibt keinen eigenständigen
\path{physics="thermomechanics"}-TOML-Modus.

Schema 2 braucht positive Dichte und Wärmekapazität. Ein Preset wie
\path{aluminium_6060_t6_20c} darf nicht gleichzeitig mit individuellen
Stoffwerten vermischt werden. Eigene Werte verwenden Leitfähigkeit in
W/(m K), Dichte in kg/m$^3$, Wärmekapazität in J/(kg K). \path{time.values}
sind streng steigende Sekunden; \path{initial_temperature} ist Kelvin.
\path{boundary.temperature} ersetzt den stationären Dirichlet-Eintrag.
\path{kind="csv"} verlangt eine Zeitspalte und explizite Wertspalte mit
SI-Einheit; die CSV-Auflösung erfolgt nur beim Dateiladen innerhalb des
Modellverzeichnisses. Eingebettete CSVs sind synthetische Lastkurven.
\path{linear} bzw. \path{previous} beschreiben Interpolation;
Extrapolation ist standardmäßig ein Fehler und \path{hold} muss explizit sein.
Profilereignisse werden ins Raster aufgenommen. Quellen/Flüsse wirken am
Intervallende linksseitig, vorgegebene Temperaturen rechtsstetig.
Diese Ereignisbehandlung ersetzt keinen Zeitkonvergenznachweis.

Schema 3 benötigt Dichte in kg/m$^3$. Modal gelten homogene Randstörungen
und \path{analysis.count}; Frequenzen werden in Hz ausgegeben, Kreisfrequenzen
in rad/s. Transient werden Zeiten in s und optional kinematische Daten
vorgegeben; elektrische Anfangswerte sind kein unabhängiger Massenzustand.
Harmonisch sind Frequenzen nichtnegativ und strikt steigend; komplexe
Lastfaktoren bestehen aus \path{real}/\path{imag}. Allgemeine Lastvektoren und
bewegte Ränder stehen zusätzlich über die direkten Solver zur Verfügung.

\section{Materialzuordnung und SI-Konventionen}
Jede aktive Volumenzelle benötigt pro Feld genau eine Materialzuordnung.
\path{regions.T}, \path{regions.u}, \path{regions.phi} bzw.
\path{regions.piezo} ordnen Physical-Groups den jeweiligen Katalognamen zu;
\path{regions.expansion} benennt bei Kopplungen den Ausdehnungsdatensatz.
Fehlende Zellen oder widersprüchliche Überlappungen sind Eingabefehler.
Ein Materialwechsel erzeugt keine automatische neue Grenzflächenbedingung.

Direkte Python-Materialobjekte erwarten SI. Model-Schema 1 besitzt nur die
explizit dokumentierten Umrechnungen, etwa GPa oder absolute \path{degC}.
\path{degC} addiert 273.15; Temperaturdifferenzen dürfen nicht mit diesem
Offset behandelt werden. \path{1/delta_degC} entspricht \path{1/K},
\path{1/degC} ist unzulässig. Relative Permittivität ist dimensionslos und
wird nur ausdrücklich als \path{relative} mit dem Schemawert von
$\varepsilon_0$ in F/m umgerechnet. Absolute Permittivität erhält die Einheit
F/m. Ein Zahlenwert in pC/N ist kein Stress-Charge-Koeffizient in C/m$^2$.

Piezo verwendet standardmäßig Stress-Charge: $C^E$, $e$, $\kappa^S$.
Strain-Charge-Daten werden vor dem Solve vollständig umgerechnet; die
Koeffizienten dürfen nicht zwischen Formen gemischt werden. Die mechanische
Voigt-Reihenfolge ist $xx,yy,zz,yz,xz,xy$ mit Engineering-Schubdehnungen.
\path{axes} enthält die lokalen Materialachsen als Spalten in globalen
Koordinaten; die dritte lokale Achse ist die übliche Polungsachse.
Provenienz erfasst Eingabewerte, Konversion und effektive globale Tensoren.
Synthetische Werte der Rezepte sind keine Herstellerempfehlung.

\section{Randwerte, Elektroden und Preflight}
\path{boundary.dirichlet} wählt entweder nullbasierte Netzknoten oder einen
benannten Flächensatz. Für \path{u} sind Teilkomponenten \path{x,y,z}
möglich; ein voller mechanischer Knoten besitzt drei DOFs. Feldindizes
nicht mit geometrischen Knotenindizes verwechseln. Direkte Solver verwenden
\path{thermal_dofs}, \path{mechanical_dofs}, \path{electric_dofs} bzw.
\path{piezoelectric_dofs} zur korrekten Abbildung.
Positiver Wärmefluss zeigt nach außen. Traktion ist N/m$^2$, Volumenkraft
N/m$^3$, Oberflächenladung C/m$^2$, Volumenladung C/m$^3$.

Eine Elektrode mit festem Potential ist äquipotentiell. Eine schwebende
Elektrode mit \path{charge=0} setzt die \emph{Gesamtladung} null;
sie ist weder punktweise $D\cdot n=0$ noch eine elektrische Referenz.
Ohne Elektrodenbeschreibung ist eine isolierte Fläche kein idealer Leiter.
Überlappende Elektroden, widersprüchliche Potentiale oder zusätzliche
Flächenladung auf derselben Elektrodenfläche werden abgelehnt.

Bereits die Konstruktion von \path{Model} prüft Schema, Materialien,
Geometriezuordnung und feldweise Verankerung. Jede zusammenhängende
stationäre Wärme-Komponente braucht Temperaturvorgabe oder wirksame
Konvektion; Mechanik braucht die sechs Starrkörpermoden verankert;
Potential braucht eine feste Gauge. Transiente Wärme kann durch positive
Kapazität ohne stationäre Temperaturreferenz lösbar sein. Preflight setzt
keine Erdung und korrigiert keine Randbedingungen. Er prüft Struktur,
keine vollständige Lastkompatibilität oder Kondition. Erst der Solve und
physikalische Nachweis beantworten die weiteren Fragen.
